%% chapter4.tex: Experimental Evaluation

\chapter{Experimental Evaluation}
\label{chap:evaluation}

\section{Introduction}

This chapter presents the experimental evaluation of the routing system
described in Chapter~\ref{chap:implementation}. The evaluation was carried
out on physical hardware in a three-node testbed that reproduces the target
operational scenario: a mobile robot transmitting a live video stream to a
stationary base station, with an intermediate relay node available to
maintain connectivity after a direct link failure.

The chapter is organised as follows. Section~\ref{sec:setup} describes the
hardware platform, the per-node network and software configuration, the video
pipeline, and the link-failure methodology. Section~\ref{sec:metrics} defines
the performance metrics and their measurement formulas.
Section~\ref{sec:experiments} describes the four experiments performed.
Section~\ref{sec:results} presents and analyses the results for each
experiment, including inter-arrival time histograms that characterise the
TDMA burst structure under direct and relay conditions.
Section~\ref{sec:comparison} discusses the comparison with the Layer~2
routing approach of Morais~\cite{morais2024synchronization}.
Section~\ref{sec:eval-summary} summarises the findings.


%% -----------------------------------------------------------------------
\section{Experimental Setup}
\label{sec:setup}

\subsection{Hardware Platform}
\label{subsec:hardware}

The testbed comprised three nodes, each assigned a fixed role throughout
all experiments. Table~\ref{tab:hardware} lists the hardware components of
each node.

\begin{table}[h]
\centering
\caption{Hardware components of each testbed node}
\label{tab:hardware}
\begin{tabular}{llll}
\toprule
\textbf{Component} & \textbf{N1 (Robot)} & \textbf{N2 (Relay)} & \textbf{N3 (Base Station)} \\
\midrule
Platform        & AlphaBot2-Pi    & Raspberry Pi 4B & Laptop (Ubuntu 22.04) \\
Processor       & RPi 4B, 4\,GB   & RPi 4B, 4\,GB   & Intel Core i7 \\
Wireless        & USB Wi-Fi (ad-hoc) & USB Wi-Fi (ad-hoc) & USB Wi-Fi (ad-hoc) \\
Camera          & Pi Camera v2    & ---             & --- \\
Motor driver    & TB6612FNG       & ---             & --- \\
Servo controller & PCA9685 (I\textsuperscript{2}C) & --- & --- \\
Gamepad         & ---             & ---             & DualShock 4 (USB) \\
\bottomrule
\end{tabular}
\end{table}

\textbf{Node 1 (N1) --- Robot.} N1 is an AlphaBot2-Pi differential-drive
robot equipped with a Raspberry Pi 4 Model B (4\,GB RAM). A Raspberry Pi
Camera Module~v2 is mounted at the front of the chassis for live video
capture. A USB Wi-Fi adapter provides wireless connectivity in IEEE~802.11
ad-hoc (IBSS) mode. The robot's two DC motors are driven by a TB6612FNG
H-bridge motor driver controlled via PWM from the Raspberry Pi GPIO header.
A PCA9685 16-channel PWM controller on the I\textsuperscript{2}C bus drives
the pan servo for camera orientation. N1 runs the RA-TDMAs+ node software
and the \texttt{alphabot\_node.py} teleoperation server, which manages video
streaming, motor control, and telemetry.

\textbf{Node 2 (N2) --- Relay.} N2 is a Raspberry Pi 4 Model B (4\,GB RAM)
with a USB Wi-Fi adapter. It acts as a stationary relay node and runs only
the RA-TDMAs+ node software; it does not originate or consume application
data.

\textbf{Node 3 (N3) --- Base Station.} N3 is a laptop running Ubuntu~22.04,
connected to the mesh via a USB Wi-Fi adapter. N3 runs the RA-TDMAs+ node
software and the \texttt{base\_station.py} operator interface, which displays
the incoming video stream and translates DualShock~4 gamepad inputs into
teleoperation commands.

All three nodes were placed on a laboratory bench within mutual radio range
so that any pair could communicate directly without a relay. This placement
allows link failures to be induced by software rather than by physical
movement, ensuring full reproducibility.

\subsection{AlphaBot2-Pi Platform}
\label{subsec:alphabot}

The AlphaBot2-Pi is an educational differential-drive robot designed for use
with a Raspberry Pi single-board computer. Its mechanical chassis carries two
independently driven wheels and a front-mounted camera gimbal. The Raspberry
Pi 4 Model B provides four ARM Cortex-A72 cores at 1.8\,GHz, 4\,GB of
LPDDR4 RAM, a 40-pin GPIO header, and an I\textsuperscript{2}C bus, all of
which are exploited by this work.

Motor control uses the TB6612FNG dual H-bridge driver. Each motor channel
accepts a PWM signal on the enable pin and two direction signals on the input
pins. \texttt{alphabot\_node.py} drives these signals at 500\,Hz via the
\texttt{RPi.GPIO} library, with duty cycle proportional to the commanded
speed in the range $[-1.0, +1.0]$. The PCA9685 servo controller is accessed
over I\textsuperscript{2}C at address \texttt{0x40} and configured for a
50\,Hz PWM frequency. Channel~0 drives the pan servo, with a calibrated
centre position of 100° corresponding to a pulse width of approximately
1.5\,ms. Infrared line sensors (left and right) and an HC-SR04 ultrasonic
distance sensor are also read by the teleoperation server and reported to
N3 as telemetry at 5\,Hz.

The platform was chosen because it matches the hardware used in the prior
work by Morais~\cite{morais2024synchronization}, enabling a direct comparison
of the two routing approaches under identical physical conditions.

\subsection{Network and Software Configuration}
\label{subsec:netconfig}

Each node was configured individually before each experiment session. The
configuration steps are grouped below by node.

\paragraph{Common configuration (all nodes).}
Each node set its USB Wi-Fi adapter to IEEE~802.11 ad-hoc (IBSS) mode on the
same SSID and channel, assigned a physical address in the
\texttt{172.20.10.0/24} subnet, and brought the interface up. A TUN virtual
interface was created and assigned a mesh address in the \texttt{10.0.0.0/24}
subnet. IP forwarding and reverse-path filter settings were applied as
described in Section~\ref{sec:kernel-config}. Table~\ref{tab:netconfig}
summarises the address assignments.

\begin{table}[h]
\centering
\caption{Network address assignments}
\label{tab:netconfig}
\begin{tabular}{lccc}
\toprule
\textbf{Interface} & \textbf{N1 (Robot)} & \textbf{N2 (Relay)} & \textbf{N3 (Base Station)} \\
\midrule
\texttt{wlan0} (physical) & \texttt{172.20.10.1/28} & \texttt{172.20.10.2/28} & \texttt{172.20.10.3/28} \\
\texttt{tun\textit{id}} (mesh)   & \texttt{10.0.0.1/24}    & \texttt{10.0.0.2/24}    & \texttt{10.0.0.3/24}    \\
\bottomrule
\end{tabular}
\end{table}

\paragraph{N1 (Robot).}
After the common configuration, \texttt{alphabot\_node.py} was started. The
script waits 10\,s for the TDMA node software to establish synchronisation
before launching the video pipeline.

\paragraph{N2 (Relay).}
The relay node required no application-level configuration beyond the common
network setup. Kernel IP forwarding and the \texttt{iptables} FORWARD rules
described in Section~\ref{sec:kernel-config} were applied to permit transit
packets between \texttt{tun\textit{id}} and \texttt{wlan0}.

\paragraph{N3 (Base Station).}
\texttt{base\_station.py} was started, which launches \texttt{ffplay} for
video display and begins listening for UDP video datagrams on port~5000 and
telemetry on port~9001. For traffic capture, \texttt{tcpdump -i any} was
started before each experiment run to record all packets for
post-experiment analysis.

For link-failure experiments, the following \texttt{iptables} rules were
applied on N3 at the moment of the simulated failure:

\begin{itemize}
  \item \texttt{iptables -I INPUT 1 -s 172.20.10.1 -j DROP}
  \item \texttt{iptables -I OUTPUT 1 -d 172.20.10.1 -j DROP}
\end{itemize}

These rules cause all frames from and to N1's physical address to be
silently discarded at the firewall level, making the direct link invisible
to the TDMA beacon receiver and to the TCP data-plane connection, without
any physical movement of nodes. Rules were cleared between repetitions
with \texttt{iptables -F}.

\subsection{Video Pipeline and Application Parameters}
\label{subsec:video-pipeline}

The video pipeline on N1 connects two processes through a Unix pipe.
\texttt{rpicam-vid} captures raw frames from the Pi Camera Module~v2 in
YUV~4:2:0 format and streams them to standard output. \texttt{ffmpeg} reads
this raw stream, encodes it with \texttt{libx264}, wraps the output in
MPEG-TS containers, and sends each transport-stream packet as a UDP datagram
to N3's mesh address (\texttt{10.0.0.3:5000}). Because the destination is
a TUN interface address, the datagram enters the TDMA routing layer
transparently: the TDMA framework reads it from \texttt{tun\textit{id}} and
delivers it to N3 via the TCP data plane, either directly or through N2
depending on the current routing table.

On N3, a Python proxy thread in \texttt{base\_station.py} receives each
UDP datagram on port~5000, records its arrival timestamp for measurement
purposes, and forwards it unchanged to \texttt{ffplay} listening on the
loopback port~5001. \texttt{ffplay} is configured for minimum display latency
using the \texttt{nobuffer}, \texttt{discardcorrupt}, \texttt{low\_delay},
and \texttt{framedrop} options. Table~\ref{tab:app-params} lists all
application parameters.

\begin{table}[h]
\centering
\caption{Application parameters used in all experiments}
\label{tab:app-params}
\begin{tabular}{lll}
\toprule
\textbf{Parameter} & \textbf{Value} & \textbf{Constant / source} \\
\midrule
\multicolumn{3}{l}{\textit{Video capture and encoding (N1)}} \\
Capture tool        & \texttt{rpicam-vid}              & \texttt{alphabot\_node.py} \\
Capture format      & YUV 4:2:0 raw                    & \texttt{--codec yuv420} \\
Encoder             & FFmpeg / \texttt{libx264}         & \texttt{alphabot\_node.py} \\
Resolution          & $640 \times 480$\,px              & \texttt{VIDEO\_WIDTH / HEIGHT} \\
Frame rate          & 15\,fps                           & \texttt{VIDEO\_FPS} \\
Target bitrate      & 500\,kbps                         & \texttt{VIDEO\_BITRATE} \\
Encoder preset      & \texttt{ultrafast}                & \texttt{alphabot\_node.py} \\
Encoder tune        & \texttt{zerolatency}              & \texttt{alphabot\_node.py} \\
GOP size            & 1 (all-intra, no inter-prediction) & \texttt{-g 1} \\
Container           & MPEG-TS                           & \texttt{-f mpegts} \\
MPEG-TS packet size & 1316\,bytes ($7 \times 188$)      & \texttt{pkt\_size=1316} \\
Stream protocol     & UDP                               & \texttt{alphabot\_node.py} \\
Stream destination  & \texttt{10.0.0.3:5000}           & \texttt{BASE\_IP:VIDEO\_PORT} \\
\midrule
\multicolumn{3}{l}{\textit{Video reception (N3)}} \\
Player              & \texttt{ffplay}                   & \texttt{base\_station.py} \\
Player flags        & \texttt{nobuffer}, \texttt{discardcorrupt}, \texttt{low\_delay}, \texttt{framedrop} & \texttt{base\_station.py} \\
Receive port        & UDP 5000                          & \texttt{VIDEO\_PORT} \\
\midrule
\multicolumn{3}{l}{\textit{Control channel (bidirectional)}} \\
Command port (N1)   & UDP 9000                          & \texttt{CMD\_PORT} \\
Command rate        & 20\,Hz                            & \texttt{CMD\_INTERVAL = 0.05\,s} \\
Telemetry port (N3) & UDP 9001                          & \texttt{TEL\_PORT} \\
Telemetry rate      & 5\,Hz                             & \texttt{TELEMETRY\_INTERVAL = 0.2\,s} \\
\midrule
\multicolumn{3}{l}{\textit{Motor and servo (N1)}} \\
Motor PWM frequency & 500\,Hz                           & \texttt{GPIO.PWM(..., 500)} \\
Speed modes         & 30\,\% / 55\,\% / 80\,\%         & \texttt{SPEED\_MODES} \\
Servo centre (pan)  & 100°                              & \texttt{SERVO\_PAN\_CENTER} \\
\bottomrule
\end{tabular}
\end{table}


%% -----------------------------------------------------------------------
\section{Metrics}
\label{sec:metrics}

\paragraph{Round-trip time (RTT).}
The RTT of the $k$-th echo exchange is the elapsed time from the moment N3
transmits a UDP echo request to the moment the corresponding reply from N1
is received:
\begin{equation}
\mathrm{RTT}_k = t^{\mathrm{rx}}_k - t^{\mathrm{tx}}_k
\label{eq:rtt}
\end{equation}
The mean RTT over $N$ samples is:
\begin{equation}
\overline{\mathrm{RTT}} = \frac{1}{N}\sum_{k=1}^{N}\mathrm{RTT}_k
\label{eq:rtt-mean}
\end{equation}
RTT captures the end-to-end latency of the mesh and, in a TDMA network,
reflects the slot structure: a request sent at an arbitrary moment within
the frame must wait until N1's slot before the reply can depart.

\paragraph{Packet delivery ratio (PDR).}
PDR quantifies the reliability of the control channel:
\begin{equation}
\mathrm{PDR} = \frac{N_{\mathrm{rx}}}{N_{\mathrm{tx}}} \times 100\,\%
\label{eq:pdr}
\end{equation}
where $N_{\mathrm{tx}}$ is the number of echo requests sent and
$N_{\mathrm{rx}}$ is the number of replies received within a timeout window.

\paragraph{Jitter.}
Jitter is the mean absolute deviation of consecutive RTT samples, measuring
the variability of the control-plane delay:
\begin{equation}
J = \frac{1}{N-1}\sum_{k=2}^{N}\bigl|\mathrm{RTT}_k - \mathrm{RTT}_{k-1}\bigr|
\label{eq:jitter}
\end{equation}

\paragraph{Throughput.}
Video stream throughput at N3 is computed from the total bytes received
during an observation window of duration $\Delta t$:
\begin{equation}
T = \frac{8 \times B_{\mathrm{total}}}{\Delta t} \quad \text{[kbps]}
\label{eq:throughput}
\end{equation}
where $B_{\mathrm{total}}$ is the sum of all UDP payload bytes recorded
by the \texttt{tcpdump} capture.

\paragraph{Convergence time.}
Convergence time quantifies how quickly the system recovers from a direct
link failure and re-establishes the video stream via the relay:
\begin{equation}
t_{\mathrm{conv}} = t_{\mathrm{first\,relay}} - t_{\mathrm{last\,direct}}
\label{eq:convergence}
\end{equation}
where $t_{\mathrm{last\,direct}}$ is the PCAP timestamp of the last video
packet received via the direct path and $t_{\mathrm{first\,relay}}$ is the
timestamp of the first packet received via N2.

\paragraph{Inter-arrival time (IA).}
The inter-arrival time between consecutive video packets is:
\begin{equation}
\mathrm{IA}_k = t_{k+1} - t_k
\label{eq:ia}
\end{equation}
where $t_k$ is the PCAP arrival timestamp of the $k$-th packet. The
distribution of $\mathrm{IA}_k$ reveals the burst structure imposed by
the TDMA slot schedule and allows verification that relay forwarding
preserves this structure unchanged.


%% -----------------------------------------------------------------------
\section{Experiments}
\label{sec:experiments}

\textbf{Experiment 1 --- Control-channel RTT and PDR.}
N3 sent UDP echo requests to N1's mesh address (\texttt{10.0.0.1}) at
10\,Hz for 60\,s under steady-state direct-link conditions, producing 600
samples. For each request, N1's node software replied immediately via the
same UDP socket. The RTT, PDR, minimum, maximum, and standard deviation
were computed from the 600 recorded samples. This experiment characterises
the baseline latency and reliability of the mesh and verifies that the
measured mean RTT is consistent with the TDMA slot ordering.

\textbf{Experiment 2 --- Video stream throughput.}
N1 streamed video continuously to N3 for approximately two minutes under
both direct-link and relay conditions. Two independent runs were performed
for each condition. Throughput was computed from the PCAP capture at N3
using Equation~\ref{eq:throughput} over the full capture duration. For
the relay condition, the \texttt{iptables} rules described in
Section~\ref{subsec:netconfig} were applied after the stream had been
running stably for at least 30\,s, and measurements were taken only after
the routing system had converged onto the relay path.

\textbf{Experiment 3 --- Routing convergence.}
Six link-failure events were induced across two independent test runs,
with three repetitions per run. Each repetition followed the same
procedure: the video stream was allowed to stabilise under direct-link
conditions for at least 30\,s, the \texttt{iptables} rules were applied
to simulate a link failure, and the PCAP capture was used to measure the
convergence time using Equation~\ref{eq:convergence}. Repetitions affected
by extraneous conditions were excluded from the timing statistics but
reported separately with their cause.

\textbf{Experiment 4 --- Packet inter-arrival timing.}
Inter-arrival times were extracted from the same PCAP captures used in
Experiment 2, separately for the direct-link and relay phases.
Histograms of the IA distribution were produced for both conditions and
compared to verify that relay forwarding preserves the TDMA burst
structure of the video stream.


%% -----------------------------------------------------------------------
\section{Results and Discussion}
\label{sec:results}

\subsection{Control-Channel Latency}
\label{subsec:res-rtt}

Table~\ref{tab:rtt} summarises the RTT statistics obtained from the 600
echo samples of Experiment~1.

\begin{table}[h]
\centering
\caption{Control-channel RTT statistics (N3 $\to$ N1, direct link, $n = 600$)}
\label{tab:rtt}
\begin{tabular}{lc}
\toprule
\textbf{Metric} & \textbf{Value} \\
\midrule
Mean RTT ($\overline{\mathrm{RTT}}$) & 50.33\,ms \\
Minimum RTT  & 17\,ms  \\
Maximum RTT  & 125\,ms \\
Standard deviation & 6.2\,ms \\
Jitter ($J$) & 4\,ms   \\
PDR          & 100\,\% \\
\bottomrule
\end{tabular}
\end{table}

The mean RTT of approximately 50\,ms is fully consistent with the TDMA
slot structure. In the slot ordering $[\text{N1}, \text{N2}, \text{N3}]$,
N1 transmits during the first 50\,ms of each 150\,ms frame. An echo
request sent by N3 at a uniformly random instant within the frame must
wait, on average, one full slot duration of 50\,ms before N1 can reply.
The expected mean wait is therefore:
\begin{equation}
\mathbb{E}[\text{wait}] = \frac{T_{\mathrm{slot}}}{2} + \frac{T_{\mathrm{slot}}}{2}
= T_{\mathrm{slot}} = 50\,\text{ms}
\label{eq:expected-rtt}
\end{equation}
which matches the observed mean exactly. The minimum of 17\,ms arises
when the request is sent just before N1's slot begins; the maximum of
125\,ms corresponds to a request sent just after N1's slot has closed,
requiring a wait of nearly one full frame. The jitter of 4\,ms is low and
reflects processing variability at N1 rather than any timing instability
in the TDMA schedule. The PDR of 100\% confirms that the UDP control
channel is fully reliable under these conditions.

\subsection{Video Stream Throughput}
\label{subsec:res-throughput}

Table~\ref{tab:throughput} presents the video stream throughput measured
at N3 under direct-link and relay conditions across two independent runs.

\begin{table}[h]
\centering
\caption{Video stream throughput at N3 (kbps)}
\label{tab:throughput}
\begin{tabular}{lccc}
\toprule
\textbf{Condition} & \textbf{Run 1} & \textbf{Run 2} & \textbf{Mean} \\
\midrule
Direct & 583.3 & 589.2 & 586.2 \\
Relay  & 570.1 & 580.5 & 575.3 \\
\midrule
\textbf{Difference} & \multicolumn{3}{c}{$-1.9\%$ (relay vs.\ direct)} \\
\bottomrule
\end{tabular}
\end{table}

The direct-link mean of 586\,kbps slightly exceeds the 500\,kbps encoder
target. Two factors account for this difference. First, the MPEG-TS
container adds approximately 15\% of overhead: every 188-byte
transport-stream packet carries a 4-byte header, and the ffmpeg pipeline
groups seven packets per UDP datagram (\texttt{pkt\_size=1316}), so
container overhead is unavoidable. Second, the encoder is configured with
\texttt{-g~1}, which disables inter-frame prediction; every frame is
encoded as an independent intra frame (I-frame), producing a higher and
more constant bitrate than a configuration that uses P-frames or B-frames.

Under relay conditions the mean throughput is 575\,kbps, a reduction of
1.9\% relative to the direct case. This difference is within the
run-to-run variability introduced by the H.264 encoder as scene content
changes, and no systematic degradation attributable to the relay hop
itself is observed. The relay's 50\,ms slot provides approximately
$0.05 \times 54 \times 10^3 = 2700$\,kb of theoretical capacity per frame,
of which the video stream consumes only $575 \times 0.150 = 86$\,kb per
frame --- well within the available budget.

\subsection{Routing Convergence}
\label{subsec:res-convergence}

Table~\ref{tab:convergence} presents the convergence times recorded across
all six link-failure repetitions of Experiment~3. Two repetitions were
excluded from the aggregate statistics due to extraneous conditions
unrelated to the routing mechanism.

\begin{table}[h]
\centering
\caption{Routing convergence time after simulated link failure}
\label{tab:convergence}
\begin{tabular}{cccc}
\toprule
\textbf{Run} & \textbf{Rep.} & \textbf{$t_{\mathrm{conv}}$ (ms)} & \textbf{Notes} \\
\midrule
1 & 1 & 2847 & --- \\
1 & 2 & 1801 & excluded (battery depletion on N1) \\
1 & 3 & 2847 & --- \\
\midrule
2 & 1 & 2847 & --- \\
2 & 2 & 2848 & --- \\
2 & 3 & ---  & excluded (residual \texttt{iptables} rule) \\
\midrule
\multicolumn{2}{l}{\textbf{Mean (valid runs)}} & \textbf{2847} & Std: 1\,ms \\
\bottomrule
\end{tabular}
\end{table}

The mean convergence time over the four valid repetitions is approximately
2847\,ms with a standard deviation of 1\,ms, indicating a highly
deterministic recovery process. This value corresponds to approximately
19 TDMA frames of 150\,ms each. The convergence chain after the direct
link is severed comprises the following sequential steps:

\begin{enumerate}
  \item \textbf{Beacon loss detection (approx.\ 450\,ms).}
  RA-TDMAs+ declares a link failure after three consecutive missed beacons.
  With a frame period of 150\,ms, this takes $3 \times 150 = 450$\,ms.

  \item \textbf{MST recomputation (tens of $\mu$s).}
  The PDR for the N1--N3 link drops to zero, triggering a rerun of Prim's
  algorithm over the updated quality graph. This completes in microseconds
  on the small three-node graph.

  \item \textbf{Netlink route installation (approx.\ 100\,$\mu$s).}
  Two host routes are installed in the kernel routing tables via Netlink
  sockets at approximately 50\,$\mu$s each.

  \item \textbf{TCP teardown and reconnection via relay (approx.\ 2.4\,s).}
  The TCP connection between N1 and N3 was established over the direct path.
  Once that path is severed, the kernel must wait for the TCP keep-alive
  probes to time out before the connection is declared dead, after which
  a new three-way handshake is initiated over the relay path
  (N1\,$\to$\,N2\,$\to$\,N3). The reconnection traverses two TDMA slots
  and completes within one to two frames (150--300\,ms), but the preceding
  teardown accounts for the bulk of the 2.4\,s delay.

  \item \textbf{First video packet via relay.}
  N1 resumes transmission over the new TCP connection; the first packet
  reaches N3 via N2.
\end{enumerate}

The dominant contribution to convergence time is TCP teardown and
reconnection, not routing computation or route installation. The
keep-alive parameters (\texttt{TCP\_KEEPIDLE = 5\,s},
\texttt{TCP\_KEEPINTVL = 2\,s}, \texttt{TCP\_KEEPCNT = 3}) set a
worst-case detection time of $5 + 3 \times 2 = 11$\,s; the observed
2.4\,s teardown time suggests that the kernel's retransmission timeout
triggered the connection reset before the keep-alive probes exhausted
their count. Reducing \texttt{TCP\_KEEPIDLE} and the retransmission
timeout further would shorten this component of the convergence delay.

\subsection{Packet Inter-Arrival Distribution}
\label{subsec:res-ia}

Figures~\ref{fig:ia-direct} and~\ref{fig:ia-relay} show the inter-arrival
time histograms of video packets at N3 under direct-link and relay
conditions, respectively, extracted from the PCAP captures of Experiment~4.

\begin{figure}[H]
\centering
\fbox{\parbox{0.85\textwidth}{\centering\vspace{3.5cm}
[Wireshark IA histogram --- direct link]
\vspace{3.5cm}}}
\caption{Inter-arrival time histogram of video packets at N3 under
direct-link conditions. The distribution is strongly bimodal: the left
peak ($<10$\,ms) corresponds to consecutive packets within the same TDMA
burst; the right peak ($\approx 130$--$150$\,ms) corresponds to the
inter-burst gap while N2 and N3 hold the channel.}
\label{fig:ia-direct}
\end{figure}

\begin{figure}[H]
\centering
\fbox{\parbox{0.85\textwidth}{\centering\vspace{3.5cm}
[Wireshark IA histogram --- relay]
\vspace{3.5cm}}}
\caption{Inter-arrival time histogram of video packets at N3 under relay
conditions. The bimodal structure is preserved identically, confirming
that N2's kernel IP forwarding introduces no additional inter-packet delay
at the burst level.}
\label{fig:ia-relay}
\end{figure}

Table~\ref{tab:ia} presents the summary statistics for both conditions.

\begin{table}[h]
\centering
\caption{Video inter-arrival time statistics at N3}
\label{tab:ia}
\begin{tabular}{lcc}
\toprule
\textbf{Metric} & \textbf{Direct} & \textbf{Relay} \\
\midrule
Mean IA (ms)    & 17.43  & 16.33  \\
Median IA (ms)  & 0.71   & 0.43   \\
Maximum IA (ms) & 149.90 & 154.96 \\
P95 IA (ms)     & 144.43 & 144.78 \\
\bottomrule
\end{tabular}
\end{table}

The distribution is strongly bimodal in both conditions. Approximately
88\% of inter-arrival intervals fall below 10\,ms, corresponding to
packets within the same TDMA transmission burst. The remaining 11\% fall
in the 130--150\,ms range, corresponding to the inter-burst gap during
which N2 and N3 hold the channel and N1 cannot transmit. This structure
arises directly from the TDMA schedule: N1 accumulates data for
approximately 100\,ms, then transmits a burst of 8--9 packets during its
50\,ms slot. At 580\,kbps with a MPEG-TS packet size of 1316 bytes,
each frame's worth of data contains:
\begin{equation}
n_{\mathrm{pkts}} = \frac{580{,}000 \times 0.150}{8 \times 1316} \approx 8.3
\label{eq:pkts-per-slot}
\end{equation}
packets per slot, consistent with the observed burst size.

The relay histogram is statistically indistinguishable from the direct
histogram. The median, P95, and maximum inter-arrival values differ by at
most 5\,ms between the two conditions. The relay hop shifts the absolute
arrival time of each burst by approximately one TDMA slot (50\,ms), but
leaves the intra-burst and inter-burst timing structure unchanged. This
confirms that the kernel IP forwarding engine at N2 forwards packets
within the same slot in which it receives them, adding no observable
jitter at the burst level.


%% -----------------------------------------------------------------------
\section{Comparison with Layer~2 Routing}
\label{sec:comparison}

A direct quantitative comparison between the Layer~3 routing system
developed in this work and the Layer~2 approach of
Morais~\cite{morais2024synchronization} is planned as part of the
experimental evaluation but was not completed within the scope of this
dissertation. The comparison would use the same three-node testbed and
identical experimental procedures (RTT, throughput, convergence, and
inter-arrival timing) applied to a clean-slate reimplementation of the
Layer~2 ARP-manipulation architecture described in
Chapter~\ref{chap:background}. The architectural differences between the
two approaches lead to the following predicted differences in experimental
behaviour:

\begin{itemize}
  \item \textbf{Convergence time}: The Layer~2 approach is expected to
  converge faster after a link failure because the TCP connection
  established at the application level survives route changes transparently;
  ARP table updates at the relay take effect within a single TDMA frame
  without requiring a TCP reconnection. The Layer~3 system must tear down
  and re-establish the TCP connection, contributing approximately 2.4\,s
  to the observed convergence time.

  \item \textbf{Throughput and inter-arrival timing}: Both approaches
  operate within the same TDMA slot constraints and are expected to produce
  equivalent throughput and burst-timing characteristics.

  \item \textbf{Observability}: The Layer~3 system is verifiable with
  standard tools (\texttt{ping}, \texttt{traceroute}, \texttt{ip route})
  that reflect the actual relay path, whereas the Layer~2 approach makes
  the relay path invisible at the IP layer.

  \item \textbf{Stability}: The Layer~2 implementation described by
  Morais exhibited runtime failures of undetermined origin. A
  reimplementation may or may not reproduce this behaviour.
\end{itemize}

Completing this comparison and quantifying these differences experimentally
is identified as the primary direction for future work.


%% -----------------------------------------------------------------------
\section{Summary}
\label{sec:eval-summary}

This chapter evaluated the routing system across four experiments on
physical hardware. The control channel achieves a mean RTT of 50\,ms
consistent with the TDMA slot structure and 100\% PDR under steady-state
conditions. Video throughput under relay conditions is within 2\% of the
direct-link baseline, confirming that the relay slot provides sufficient
capacity for the target stream at 500\,kbps. Routing convergence after a
link failure averages 2847\,ms, a highly deterministic value dominated by
TCP connection teardown and re-establishment rather than by routing
computation or route installation, both of which complete in microseconds.
The inter-arrival timing analysis confirms that kernel IP forwarding at N2
preserves the TDMA burst structure of the video stream exactly, adding no
observable jitter at the burst level. A quantitative comparison with the
Layer~2 approach of Morais~\cite{morais2024synchronization} is identified as
the main direction for future work.
